\documentclass[11pt,letterpaper]{article}
% Author: Robert Warren

\usepackage{iftex}
\ifPDFTeX
  \usepackage[utf8]{inputenc}
  \usepackage[T1]{fontenc}
  \usepackage{lmodern}
\fi
\usepackage{w3c-recommendation}

\WThreeCSetup{
  title={RDF 1.2 Concepts and Abstract Data Model},
  subtitle={W3C Recommendation-style LaTeX document},
  status={Working Draft},
  date={11 September 2026},
  version={https://www.w3.org/TR/rdf12-concepts/},
  latest={https://www.w3.org/TR/rdf12-concepts/},
  editors={Example Editor (Example Organization)},
  feedback={\href{https://github.com/w3c/}{GitHub issue tracker}},
  organization={World Wide Web Consortium},
  documentnumber={W3C Example 2026/09},
  copyright={Copyright © 2026 World Wide Web Consortium.}
}

\begin{document}
\WThreeCStart
\maketitle
\WThreeCMoreDetails{%
  This example demonstrates the document metadata and navigation structure
  used by W3C technical recommendations. The HTML build places the table of
  contents in a fixed left rail, while the PDF build keeps the same content in
  ordinary front matter.
}
\WThreeCContents

\WThreeCAbstract{%
  This example presents a reusable LaTeX layout for standards documents. It
  provides a title and metadata masthead, an abstract, a status section,
  readable section typography, and a generated table of contents.
}

\WThreeCStatus{%
  This section describes the status of this document at the time of its
  publication. It is included as normal document text so authors can replace
  it with the applicable W3C publication language.
}

\section{Introduction}

This section demonstrates the main readable content column. Links use the same
blue treatment as the W3C recommendation layout.

\subsection{Document structure}

The package keeps the document source semantic: use ordinary \LaTeX{}
\verb|\section|, \verb|\subsection|, figures, tables, and bibliographies.

\section{Conformance}

The PDF version uses a conventional table of contents after the title block.
The make4ht version moves that same table of contents into the left navigation
rail and adds responsive CSS for narrower screens.

\subsection{Implementation notes}

The style does not require external images. The default masthead uses a text
W3C mark; an image path can be supplied with the \verb|logo| setup option when
an organization logo is available.

\WThreeCEnd
\end{document}
